\section{Data Augmentation}
\paragraph{Purpose}
The best way to make a model generalize better is to train it on more data. When data is
limited, \tB{create fake data by transforming the existing training examples}.

\paragraph{Theory}
Apply to each training input $\bm{x}$ a transformation $T$ that \uB{does not change its
label} $y$, and add the pair $\left(T(\bm{x}), y\right)$ to the training set.
\begin{itemize}
    \item \textbf{Images}: translations, flips, crops, small rotations, color jitter.
    \item \textbf{Text}: synonym replacement, back-translation (translate to another
        language and back).
    \item \textbf{Any input}: injecting random noise (see noise robustness).
\end{itemize}

\paragraph{Weaknesses}
The transformation must be \tB{label-preserving for the task}: rotating a ``6'' by
$180^{\circ}$ gives a ``9'', and flipping ``b'' gives ``d''. Augmentation is applied
to the training set only, never to the validation or test sets.
